% ID: FIS-TER-PAS-001
% Titolo: Acqua spruzzata su oro fuso
% Disciplina: Fisica
% Area: Termologia
% Argomento: Passaggi di stato
% Sottoargomento: Vaporizzazione e solidificazione
% Classe: Seconda liceo scientifico
% Difficolta: 4
% Tipo: problema_numerico
% Risultato: circa 4,5 g di acqua
% Tempo_stimato: 12 min
% Competenze: bilancio termico, calore latente, interpretazione dei passaggi di stato
% Prerequisiti: calore specifico, calore latente, temperatura di ebollizione
% Tag: termologia, passaggi_stato, calore_latente, vaporizzazione, oro, acqua
% Fonte: verifica_fisica_termologia_anonimizzata
% Autore: Riccardo Carli
% Licenza: CC-BY-SA-4.0

\begin{esercizio}
Su \(180\,\mathrm{g}\) di oro fuso a \(1063\,^\circ\mathrm{C}\) viene
spruzzata acqua a \(23\,^\circ\mathrm{C}\). L'acqua si trasforma in vapore a
\(100\,^\circ\mathrm{C}\) e l'oro solidifica senza cambiare temperatura.

Calcola la quantita di acqua spruzzata sull'oro.
\end{esercizio}

\begin{soluzione}
L'oro si trova alla temperatura di fusione e solidifica senza raffreddarsi.
Il calore ceduto dall'oro e quindi il suo calore latente di solidificazione:
\[
Q_\text{oro}=m_\text{oro}L_f.
\]
Assumendo per l'oro \(L_f\simeq 6{,}45\cdot 10^4\,\mathrm{J/kg}\):
\[
Q_\text{oro}=0{,}180\cdot 6{,}45\cdot 10^4
\simeq 1{,}16\cdot 10^4\,\mathrm{J}.
\]

L'acqua assorbe calore prima per riscaldarsi da
\(23\,^\circ\mathrm{C}\) a \(100\,^\circ\mathrm{C}\), poi per vaporizzare:
\[
Q_\text{acqua}=m\left[c_\text{acqua}(100-23)+L_v\right].
\]
Con \(c_\text{acqua}\simeq 4186\,\mathrm{J/(kg\,K)}\) e
\(L_v\simeq 2{,}26\cdot 10^6\,\mathrm{J/kg}\):
\[
Q_\text{acqua}=m\left[4186\cdot 77+2{,}26\cdot 10^6\right].
\]

Ponendo \(Q_\text{acqua}=Q_\text{oro}\):
\[
m=\frac{1{,}16\cdot 10^4}{4186\cdot 77+2{,}26\cdot 10^6}
\simeq 4{,}5\cdot 10^{-3}\,\mathrm{kg}.
\]

La quantita di acqua spruzzata e quindi circa:
\[
m\simeq 4{,}5\,\mathrm{g}.
\]
\end{soluzione}
